\documentclass[a4paper,11pt]{article}
\usepackage[utf8]{inputenc}
\usepackage[T1]{fontenc}
\usepackage{lmodern}
\usepackage{amsmath,amssymb,amsthm}
\usepackage{mathtools}
\usepackage{booktabs}
\usepackage{longtable}
\usepackage{array}
\usepackage{hyperref}
\usepackage{graphicx}
\usepackage{float}
\usepackage{geometry}
\geometry{margin=25mm}

\hypersetup{
  colorlinks=true,
  linkcolor=blue,
  urlcolor=blue,
  citecolor=blue
}

\theoremstyle{definition}
\newtheorem{theorem}{Theorem}[section]
\newtheorem{proposition}[theorem]{Proposition}
\newtheorem{lemma}[theorem]{Lemma}
\newtheorem{corollary}[theorem]{Corollary}
\newtheorem{definition}[theorem]{Definition}
\newtheorem{remark}[theorem]{Remark}
\newtheorem{observation}[theorem]{Observation}

\setlength{\parindent}{1em}
\setlength{\parskip}{0.5em}

\providecommand{\tightlist}{\setlength{\itemsep}{0pt}\setlength{\parskip}{0pt}}
\begin{document}

\section{\texorpdfstring{Dimensional Interpretation of Geodesic
Structure at the \(R \to 0\)
Limit}{Dimensional Interpretation of Geodesic Structure at the R \textbackslash to 0 Limit}}\label{dimensional-interpretation-of-geodesic-structure-at-the-r-to-0-limit}

\textbf{--- An Integrated Description of Nested Structure and
Dimensional Hierarchy ---}

\begin{center}\rule{0.5\linewidth}{0.5pt}\end{center}

\textbf{Author}: Noriaki Kihara (WF System Co., Ltd.)

\textbf{Date}: April 2026

\textbf{DOI:}
\href{https://doi.org/10.5281/zenodo.19538106}{10.5281/zenodo.19538106}

\begin{center}\rule{0.5\linewidth}{0.5pt}\end{center}

\subsection{Abstract}\label{abstract}

In this paper, based on the central projection framework established in
preceding papers {[}1{]}--{[}9{]}, we provide an integrated description
of the geodesic structure at the \(R \to 0\) limit, unifying three
perspectives: the nested structure from Paper {[}4{]}, the dimensional
hierarchy from Paper {[}6{]}, and the Schwarzschild--de Sitter exact
solution from Paper {[}9{]}.

Paper {[}7{]} organized the bridging from the geodesic of a parent
subjective space to a child subjective space at an arbitrary point.
Paper {[}9{]} constructed the Schwarzschild--de Sitter metric as a
spherically symmetric static exact solution of the vacuum Einstein
equation from Paper {[}1{]}, and described the disappearance of the
horizon structure at the \(R \to 0\) limit. Building on these results,
the present paper focuses specifically on the \(R \to 0\) limit and
describes the following observations. Namely, the structure exhibited by
a subjective space at the \(R \to 0\) limit corresponds geometrically to
the description of a zero-dimensional subjective space in Paper {[}6{]},
and the nested structure of Paper {[}4{]} Proposition 4.1 confirms that
from this limit point, a new subjective space with any curvature radius
\(R' > 0\) can be constructed.

\textbf{Explicit statement}: This paper proves nothing new. It merely
organizes the geometric meaning of the \(R \to 0\) limit by combining
the nested structure of Paper {[}4{]}, the dimensional hierarchy of
Paper {[}6{]}, the geodesic bridging of Paper {[}7{]}, and the exact
solution of Paper {[}9{]}. Correspondence with concepts in physics such
as geodesic completeness, singularity theorems, and black hole
singularities is outside the scope of this paper.

\begin{center}\rule{0.5\linewidth}{0.5pt}\end{center}

\subsection{1. Introduction}\label{introduction}

Preceding Paper {[}4{]} demonstrated geometric regularity at both limits
of the curvature radius \(R\), pointed out the structural tension at the
\(R \to 0\) limit, and proved the existence of the nested structure in
Proposition 4.1. Paper {[}6{]} organized the geometric properties of
subjective spaces from zero to four dimensions. Paper {[}7{]} described
the bridging between geodesics of parent and child subjective spaces.

The purpose of this paper is to integrate these three results and
provide a description focused on the geodesic structure at the
\(R \to 0\) limit. Specifically, we address the following question.

\textbf{Question}: When tracing a geodesic on a parent subjective space
\(\Pi_R\) in the direction \(R \to 0\), which stage of the dimensional
hierarchy in Paper {[}6{]} does the structure appearing at the limit
correspond to, and what geometric degrees of freedom remain from there?

This paper proves nothing new. It merely describes combinations of
existing propositions.

\begin{center}\rule{0.5\linewidth}{0.5pt}\end{center}

\subsection{2. Preliminaries}\label{preliminaries}

\subsubsection{2.1 Terminology and Summary of Preceding
Results}\label{terminology-and-summary-of-preceding-results}

We restate the results from preceding papers used in this paper.

\textbf{Paper {[}4{]} Proposition 4.1 (Existence of nested structure)}:
For any point \(p\) in the interior of any subjective space \(\Pi_R\), a
subjective space \(\Pi_{R'}\) via a new central projection centered at
\(p\) exists for any \(R' > 0\). \(\Pi_{R'}\) possesses the same
geometric structure as the construction in Paper {[}1{]}.

\textbf{Paper {[}4{]} §3.2 (Structural tension)}: At the \(R \to 0\)
limit, the upper bound of geodesic deviation \(|\xi|^2_{\max}(R)\)
approaches zero, while the local curvature diverges at order \(1/R^2\).

\textbf{Paper {[}6{]} §2 (Zero-dimensional subjective space)}: A
zero-dimensional subjective space is a world that distinguishes only
``exists / does not exist,'' possesses no internal structure, has no
geodesic deviation, and there exists no means from within to perceive
whether it is governed by a curvature radius.

\textbf{Paper {[}7{]} §4 (Geodesic bridging)}: At any point \(p\) on a
geodesic \(\gamma\) in the parent subjective space \(\Pi_R\), a child
subjective space \(\Pi_{R'}\) can be constructed with \(p\) as the
launch point, and geodesics on \(\Pi_{R'}\) are independently
well-defined.

\subsubsection{2.2 Positioning of This
Paper}\label{positioning-of-this-paper}

Paper {[}7{]} described, in general terms, the bridging from a geodesic
of a parent subjective space to a child subjective space at \textbf{any
point}. This paper applies that general description specifically to the
\textbf{limit} \(R \to 0\) and, by adding the correspondence with the
dimensional hierarchy of Paper {[}6{]}, organizes the geometric meaning
of the limit.

\begin{center}\rule{0.5\linewidth}{0.5pt}\end{center}

\subsection{\texorpdfstring{3. Correspondence Between the \(R \to 0\)
Limit and Zero-Dimensional Subjective
Space}{3. Correspondence Between the R \textbackslash to 0 Limit and Zero-Dimensional Subjective Space}}\label{correspondence-between-the-r-to-0-limit-and-zero-dimensional-subjective-space}

\subsubsection{3.1 Vanishing of Geodesic Deviation and Dimensional
Interpretation}\label{vanishing-of-geodesic-deviation-and-dimensional-interpretation}

By Paper {[}4{]} §3.2, at the \(R \to 0\) limit, the upper bound of
geodesic deviation \(|\xi|^2_{\max}(R) \to 0\). That is, the deviation
measurable by an internal observer approaches zero at this limit.

On the other hand, in the description of the zero-dimensional subjective
space in Paper {[}6{]} §2, geodesic deviation is stated to ``not
exist.''

These two descriptions arrive at the same structural feature via
different paths. The former, as the \(R \to 0\) limit of an
\(n\)-dimensional subjective space, and the latter, as an intrinsic
property of a zero-dimensional subjective space, both describe the state
in which ``geodesic deviation is unobservable.''

\subsubsection{3.2 Disappearance of Internal Structure and Dimensional
Interpretation}\label{disappearance-of-internal-structure-and-dimensional-interpretation}

In Paper {[}6{]} §2, the internal structure of a zero-dimensional
subjective space is stated to ``not exist,'' and there exists no means
from within to perceive whether it is governed by a curvature radius
\(R\).

At the \(R \to 0\) limit, the subjective space \(\Pi_R\) ``shrinks to a
point'' (Paper {[}4{]} Question 3.2). An internal observer in a
subjective space that has shrunk to a point can distinguish neither
direction nor distance. This is geometrically consistent with the
description of the zero-dimensional subjective space in Paper {[}6{]}
§2.

\subsubsection{3.3 Description of the
Correspondence}\label{description-of-the-correspondence}

Summarizing the above, we obtain the following observation.

\textbf{Observation 3.1}: The \(R \to 0\) limit of an \(n\)-dimensional
subjective space \(\Pi_R\) corresponds geometrically to the description
of a zero-dimensional subjective space in Paper {[}6{]} in the following
respects.

\begin{enumerate}
\def\labelenumi{\arabic{enumi}.}
\tightlist
\item
  The measurable value of geodesic deviation approaching zero (Paper
  {[}4{]} §3.2) corresponds to the non-existence of geodesic deviation
  in a zero-dimensional subjective space (Paper {[}6{]} §2).
\item
  The subjective space shrinking to a point (Paper {[}4{]} Question 3.2)
  corresponds to a zero-dimensional subjective space possessing no
  internal structure (Paper {[}6{]} §2).
\item
  The internal observer becoming unable to perceive the curvature radius
  corresponds to the non-existence of any means to perceive the
  curvature radius in a zero-dimensional subjective space (Paper {[}6{]}
  §2).
\end{enumerate}

This observation contains no new proof. It is a juxtaposition of
existing descriptions from Papers {[}4{]} and {[}6{]}.

\begin{center}\rule{0.5\linewidth}{0.5pt}\end{center}

\subsection{4. Launching a Child Subjective Space from the Limit
Point}\label{launching-a-child-subjective-space-from-the-limit-point}

\subsubsection{4.1 Application of the Nested
Structure}\label{application-of-the-nested-structure}

By Paper {[}4{]} Proposition 4.1, from any point \(p\) in the interior
of the subjective space \(\Pi_R\), a child subjective space \(\Pi_{R'}\)
can be launched. As long as \(R > 0\), \(p\) is an interior point of
\(\Pi_R\) and satisfies the premises of Proposition 4.1.

At the \(R \to 0\) limit, \(\Pi_R\) approaches a zero-dimensional
structure as described in §3. However, Proposition 4.1 holds for any
value of \(R > 0\), and the curvature radius \(R'\) of the launched
child subjective space \(\Pi_{R'}\) can take any positive value
independently of \(R\).

\subsubsection{4.2 Connection with the Dimensional
Hierarchy}\label{connection-with-the-dimensional-hierarchy}

Combining Observation 3.1 from §3 with Proposition 4.1, we obtain the
following description.

\textbf{Observation 4.1}: In the process where the parent subjective
space \(\Pi_R\) approaches a zero-dimensional structure as \(R \to 0\),
from any point in the asymptotic region, a child subjective space
\(\Pi_{R'}\) with any \(R' > 0\) can be constructed by Paper {[}4{]}
Proposition 4.1. The child subjective space \(\Pi_{R'}\) possesses the
same geometric structure as the construction in Paper {[}1{]}, and
within it, geometric properties of the corresponding dimension unfold
according to the description in Paper {[}6{]} depending on the value of
\(R'\).

Rephrasing this observation in the context of the dimensional hierarchy
of Paper {[}6{]}:

\begin{itemize}
\tightlist
\item
  The zero-dimensional limit of the parent subjective space corresponds
  to the ``exists / does not exist'' world of Paper {[}6{]} §2.
\item
  However, from the ``exists'' side of that world, by Paper {[}4{]}
  Proposition 4.1, there exist geometric degrees of freedom for a new
  subjective space to be launched with any \(R' > 0\).
\item
  The launched child subjective space \(\Pi_{R'}\) possesses dimensional
  structure from Paper {[}6{]} §3 onward, depending on the value of
  \(R'\).
\end{itemize}

\subsubsection{4.3 Scope of Claims in This
Section}\label{scope-of-claims-in-this-section}

What has been stated in this section is limited to the following.

\begin{enumerate}
\def\labelenumi{\arabic{enumi}.}
\tightlist
\item
  In the limit region \(R \to 0\), Proposition 4.1 remains applicable as
  long as \(R > 0\).
\item
  The curvature radius \(R'\) of the launched child subjective space is
  independent of the parent's \(R\).
\item
  The combination of these allows a description in which the parent's
  zero-dimensional limit and the child's finite-dimensional structure
  coexist.
\end{enumerate}

This section contains no new proof. It is a combination of Proposition
4.1 and the descriptions in Paper {[}6{]}.

\begin{center}\rule{0.5\linewidth}{0.5pt}\end{center}

\subsection{5. Geodesic Tracing and Dimensional
Transition}\label{geodesic-tracing-and-dimensional-transition}

\subsubsection{\texorpdfstring{5.1 Tracing Geodesics on the Parent
Subjective Space in the \(R \to 0\)
Direction}{5.1 Tracing Geodesics on the Parent Subjective Space in the R \textbackslash to 0 Direction}}\label{tracing-geodesics-on-the-parent-subjective-space-in-the-r-to-0-direction}

As described in Paper {[}7{]} §3, when tracing a geodesic \(\gamma\) on
the parent subjective space \(\Pi_R\), the upper bound of geodesic
deviation \(|\xi|^2_{\max}(R)\) approaches zero in the \(R \to 0\)
direction.

By Observation 3.1 of §3, this asymptotic behavior can be interpreted
dimensionally as a dimensional collapse of the parent subjective space.
That is, for the internal observer of the parent subjective space, the
tracing of \(\gamma\) appears as a process in which dimensional
structure is lost.

\subsubsection{5.2 Restart of Geodesics in the Child Subjective
Space}\label{restart-of-geodesics-in-the-child-subjective-space}

By Paper {[}7{]} §4, in the child subjective space \(\Pi_{R'}\) with a
point \(p\) on \(\gamma\) as the launch point, a geodesic \(\gamma'\)
departing from \(p\) is well-defined. By Observation 4.1 of §4, this
launching is possible even in the limit region \(R \to 0\), and
\(\gamma'\) on \(\Pi_{R'}\) is defined within a dimensional structure
corresponding to \(R'\).

\subsubsection{5.3 Description as Dimensional
Transition}\label{description-as-dimensional-transition}

Combining §5.1 and §5.2, we obtain the following description.

\textbf{Observation 5.1}: The process of tracing a geodesic \(\gamma\)
on the parent subjective space \(\Pi_R\) in the \(R \to 0\) direction
can be geometrically described as the following two stages.

\textbf{Stage 1 (Dimensional collapse)}: In the parent subjective space
\(\Pi_R\), the upper bound of geodesic deviation \(|\xi|^2_{\max}(R)\)
approaches zero, and the subjective space approaches the
zero-dimensional structure of Paper {[}6{]} §2. For the internal
observer of the parent, the tracing of the geodesic appears as a loss of
dimensional structure.

\textbf{Stage 2 (Dimensional reconstruction)}: By Paper {[}4{]}
Proposition 4.1, from a point in the limit region, a child subjective
space \(\Pi_{R'}\) can be constructed with any \(R' > 0\), and within
the child subjective space, geometric structure unfolds according to the
dimensional hierarchy of Paper {[}6{]}. For the internal observer of the
child, the geodesic \(\gamma'\) is well-defined within a dimensional
structure corresponding to \(R'\).

In this two-stage description, the ``terminus'' of the geodesic in the
parent subjective space is described not as the annihilation of
geometric structure, but as a transition point of dimensional structure.
The disappearance of structure as seen by the internal observer of the
parent and the existence of structure as seen by the internal observer
of the child coexist without contradiction by Paper {[}3{]} Theorem 5.1
(relativity of observation).

\subsubsection{5.4 Scope of Claims in This
Section}\label{scope-of-claims-in-this-section-1}

What has been stated in this section is limited to the following.

\begin{enumerate}
\def\labelenumi{\arabic{enumi}.}
\tightlist
\item
  By combining the dimensional correspondence of §3 and the nested
  structure of §4, the behavior of geodesics at the \(R \to 0\) limit
  can be described as a ``dimensional transition.''
\item
  This description is consistent for observers in both the parent and
  child (by Paper {[}3{]} Theorem 5.1).
\item
  ``Dimensional transition'' is a term for geometric description and
  does not assert a physical process.
\end{enumerate}

This section contains no new proof.

\begin{center}\rule{0.5\linewidth}{0.5pt}\end{center}

\subsection{6. Discussion}\label{discussion}

In this paper, we described the geodesic structure at the \(R \to 0\)
limit in a form that integrates four elements: the nested structure of
Paper {[}4{]}, the dimensional hierarchy of Paper {[}6{]}, the geodesic
bridging of Paper {[}7{]}, and the Schwarzschild--de Sitter exact
solution of Paper {[}9{]}.

The descriptions in this paper are limited to the following.

\begin{enumerate}
\def\labelenumi{\arabic{enumi}.}
\tightlist
\item
  The \(R \to 0\) limit corresponds geometrically to the description of
  a zero-dimensional subjective space in Paper {[}6{]} (Observation
  3.1).
\item
  From the region of the zero-dimensional limit, child subjective spaces
  can be constructed, and within them dimensional structure unfolds
  (Observation 4.1).
\item
  Tracing geodesics of the parent subjective space in the \(R \to 0\)
  direction can be described as a two-stage process of dimensional
  collapse and dimensional reconstruction (Observation 5.1).
\end{enumerate}

Paper {[}9{]} supports this dimensional transition at the level of exact
solutions. Namely, the process in which the cosmological horizon
vanishes at \(R \to 0\) and the region of the internal observer is lost
(Paper {[}9{]} §4) corresponds to ``Stage 1 (Dimensional collapse)'' of
this paper, and the recovery of the horizon structure in the child
subjective space (Paper {[}9{]} §5) corresponds to ``Stage 2
(Dimensional reconstruction).''

Furthermore, regarding the stable dimension of the child subjective
space launched in Stage 2, Paper {[}8{]} provides a combinatorial
consideration. Paper {[}8{]} showed that, starting from zero-dimensional
structural information and under assumptions of discreteness and
stability, four dimensions are selected as the minimal stable dimension.
The destination of ``dimensional reconstruction'' in Observation 5.1 of
this paper connects to the discussion of Paper {[}8{]}.

The ``dimensional transition'' described in this paper is a description
within the framework of geometry and does not assert a physical process.
In particular, correspondence with established concepts in differential
geometry such as geodesic completeness, singularity theorems in general
relativity, and the behavior of geodesics in black hole spacetimes lies
outside the geometric framework and exceeds the scope of this paper.

The possible physical interpretation of the descriptions in this paper
is left to the reader's judgment.

\begin{center}\rule{0.5\linewidth}{0.5pt}\end{center}

\subsection{7. Conclusion}\label{conclusion}

In the central projection framework, the \(R \to 0\) limit of the
subjective space \(\Pi_R\) corresponds geometrically to the description
of a zero-dimensional subjective space in Paper {[}6{]}. By the nested
structure of Paper {[}4{]} Proposition 4.1, from the region of this
zero-dimensional limit, a child subjective space with any \(R' > 0\) can
be constructed, and within the child, dimensional structure unfolds.
Tracing geodesics of the parent subjective space in the \(R \to 0\)
direction is described as a two-stage process of dimensional collapse
and dimensional reconstruction, and the ``terminus'' of the geodesic is
organized not as the annihilation of geometric structure but as a
transition point of dimensional structure. All of these are described as
combinations of existing propositions from Papers {[}1{]}--{[}7{]}.
Physical interpretation is outside the scope of this paper and is left
to the reader.

\begin{center}\rule{0.5\linewidth}{0.5pt}\end{center}

\subsection{References}\label{references}

{[}1{]} Noriaki Kihara, ``Geometric Formalization of 4-Dimensional
Spacetime via Central Projection,'' Zenodo, DOI: 10.5281/zenodo.19427780
(2026).

{[}2{]} Noriaki Kihara, ``Geometric Symmetries of the Central
Projection: Mathematical Foundations of the Multi-Axis Model,'' Zenodo,
DOI: 10.5281/zenodo.19434932 (2026).

{[}3{]} Noriaki Kihara, ``Relativity of Observation in Multiple
Subjective Spaces: Geometric Consequences of Central Projection
Symmetries,'' Zenodo, DOI: 10.5281/zenodo.19435162 (2026).

{[}4{]} Noriaki Kihara, ``On the Limits of the Curvature Radius in
Central Projection --- The Cases R $\to$ $\infty$ and R $\to$ 0,'' Zenodo, DOI:
10.5281/zenodo.19526549 (2026).

{[}5{]} Noriaki Kihara, ``On Dimension Addition to Background Spaces and
Subjective Spaces,'' Zenodo, DOI: 10.5281/zenodo.19526913 (2026).

{[}6{]} Noriaki Kihara, ``Subjective Spaces from Zero to Four
Dimensions,'' Zenodo, DOI: 10.5281/zenodo.19533292 (2026).

{[}7{]} Noriaki Kihara, ``On the Continuity of Geodesics in Subjective
Spaces --- Bridging Geodesics from Parent to Child Subjective Spaces,''
Zenodo, DOI: 10.5281/zenodo.19533299 (2026).

{[}8{]} Noriaki Kihara, ``Diameter of the Circumscribed Hypersphere of a
Unit Four-Dimensional Hyperrectangle,'' Zenodo, DOI:
10.5281/zenodo.19533313 (2026).

{[}9{]} Noriaki Kihara, ``Schwarzschild--de Sitter Exact Solution in the
Central Projection Framework --- Spherically Symmetric Static Exact
Solution of the Vacuum Einstein Equation from Paper {[}1{]},'' Zenodo,
DOI: 10.5281/zenodo.19538098 (2026).

\begin{center}\rule{0.5\linewidth}{0.5pt}\end{center}

\end{document}
